\subsection{Background sources and convergence maps}
\label{sec:maps}

\subsubsection{Which galaxies carry the cluster signal?}

A calibrated shape catalog supplies samples of the lensing field at discrete sky positions. Before these samples can be combined into a cluster map, the galaxies carrying the cluster signal must be identified. A well-measured shape is useful for cluster lensing only if the galaxy lies behind the cluster. Foreground and cluster-member galaxies do not carry the same cluster-lensing signal. Under the idealized assumption that contaminants have zero mean shape, a fraction $f$ of such objects in an otherwise uniformly weighted sample reduces the mean signal by a factor $1-f$. Spatially varying contamination can then alter the inferred map morphology as well as its amplitude~\cite{schneider_weak_2006,mandelbaum_weak_2018}.

A photometric color is a difference of magnitudes in two bands, and therefore measures a flux ratio. Colors contain information about a galaxy's spectrum and redshift, although different spectra can produce similar colors. LITB~IV uses the $(u-b,b-g)$ color plane to select background candidates, calibrating the selection with external redshifts. This is separate from selecting objects with reliable shapes. The retained sources also span a redshift distribution, so converting a convergence map into physical surface density requires the corresponding lensing-efficiency distribution, through $D_{ls}/D_s$ in Eq.~\eqref{eq:convergence}~\cite{saha_lensing_2026,schneider_weak_2006}.

\subsubsection{Kaiser--Squires reconstruction}

Once the background-source sample is defined, the spatially averaged shapes estimate the reduced-shear field. Recovering convergence then uses the common potential in Eq.~\eqref{eq:potential-derivatives}. On a small, approximately planar sky patch, choose a Fourier convention for which $\partial_j$ becomes $ik_j$. For $\boldsymbol{k}\ne0$,
\begin{align}
    \widetilde\kappa&=-\tfrac12(k_1^2+k_2^2)\widetilde\psi, &
    \widetilde\gamma_1&=-\tfrac12(k_1^2-k_2^2)\widetilde\psi, &
    \widetilde\gamma_2&=-k_1k_2\widetilde\psi.
    \label{eq:fourier-potential}
\end{align}
Eliminating $\widetilde\psi$ gives the Kaiser--Squires relation
\begin{equation}
    \widetilde\kappa_E+i\widetilde\kappa_B
    =D^*(\boldsymbol{k})(\widetilde\gamma_1+i\widetilde\gamma_2),
    \qquad
    D(\boldsymbol{k})=\frac{k_1^2-k_2^2+2ik_1k_2}{k_1^2+k_2^2}.
    \label{eq:kaiser-squires}
\end{equation}
For an ideal shear field derived from a scalar lensing potential, $\kappa_E=\kappa$ and $\kappa_B=0$. The E/B notation describes the two components of this field decomposition; it does not refer to electromagnetic fields. In observations, the input is a spatial estimate from calibrated galaxy shapes, with $g\simeq\gamma$ assumed if the linear inversion is used directly~\cite{kaiser_mapping_1993,schneider_weak_2006}.

The zero Fourier mode is unconstrained: shear alone does not fix the mean convergence of the patch. More generally, for a single source plane, the transformation
\begin{equation}
    \kappa'=\lambda\kappa+(1-\lambda),\qquad
    \gamma'=\lambda\gamma
    \label{eq:mass-sheet}
\end{equation}
leaves the reduced shear unchanged. This is the mass-sheet degeneracy. Assigning a value to the mean map is a boundary or normalization choice, not an independent measurement of the missing mass sheet. Finite fields, masks, and nonuniform sampling introduce further complications because the ideal Fourier relation assumes access to the continuous field~\cite{schneider_weak_2006}.

\subsubsection{Noise, null tests, and comparison maps}

The source density limits map resolution: finer bins contain fewer galaxies and more shape noise. Smoothing suppresses noise but broadens genuine structures. Comparing maps requires consistent selections, coordinates, masks, and smoothing. Galaxy shapes must be expressed in the reconstruction's coordinate basis using their spin-2 transformation~\cite{schneider_weak_2006}.

LITB~IV estimates map noise by randomizing galaxy orientations and repeating the reconstruction. At each map position, the scatter across these realizations gives an estimate $\sigma_\kappa$, from which a local signal-to-noise map can be formed:
\begin{equation}
    \nu(\boldsymbol{\theta})=
    \frac{\kappa_E(\boldsymbol{\theta})}{\sigma_\kappa(\boldsymbol{\theta})}.
    \label{eq:map-snr}
\end{equation}
This estimates noise for the adopted catalog and reconstruction, excluding calibration and source-redshift uncertainties. Smoothing correlates neighboring pixels, so a local peak's signal-to-noise ratio is not a global detection significance~\cite{saha_lensing_2026,schneider_weak_2006}.

A B-mode map is a useful null test for the leading scalar-potential lensing model. Residual PSF patterns, noise, and boundaries can contribute to it. Conversely, a multiplicative error can rescale a pure E-mode signal without generating a B mode, so a small B mode cannot certify the E-mode calibration~\cite{schneider_weak_2006,mandelbaum_weak_2018}.

LITB~IV compares lensing maps with X-ray emission and the density of red-sequence cluster members, galaxies selected from a narrow locus in color--magnitude space. Lensing traces total projected mass; the other maps trace hot gas and member galaxies, respectively. Their structures can differ, especially in mergers~\cite{saha_lensing_2026}.
